\begin{recipe}
[% 
    preparationtime = {\unit[15]{min}},
    bakingtime={\unit[40]{min}},
    portion = {\portion{2}},
    source = {Adrian Hieber},
    calory = {620 kcal (905 kcal) | 30g (38g) Protein | 72g (127g) KH | 23g (28g) Fett},
    rating = {5},
    price = {1,54€ (2,03€)},
]
{Dal Makhani}
    
% \graph
% {% pictures
%     small=pic/dal_small,
%     big=pic/dal_big
% }

\introduction{%
    Mein Lieblingsgericht. Einfach schön gewürzig und cremig. Es ist zwar nicht das ansehnlichstes, aber dafür umso besser im Geschmack.
}

\ingredients{%
    1 & Rote Zwiebel\index{Gemüse \& Pilze!Zwiebeln}\\
    1 cm & Ingwer\index{Gemüse \& Pilze!Ingwer}\\
    3 & Knoblauchzehen\index{Gemüse \& Pilze!Knoblauch}\\
    1 & Karotte\index{Gemüse \& Pilze!Karotten}\\
    150 g & Beluga-Linsen\index{Hülsenfrüchte!Beluga-Linsen}\\
    200 ml & Passierte Tomaten\index{Vorrat \& Konserven!Tomaten (Konserve)}\\
    150 ml & Kokosmilch\index{Milchprodukte \& Alternativen!Kokosmilch}\\
    300 ml & Gemüsebrühe\index{Vorrat \& Konserven!Gemüsebrühe}\\
    1 TL & Margarine\index{Milchprodukte \& Alternativen!Margarine}\\
    2 TL & Garam Masala\index{Kräuter \& Gewürze!Garam Masala}\\
    1 TL & Paprikapulver\index{Kräuter \& Gewürze!Paprikapulver}\\
    1 TL & Kreuzkümmel\index{Kräuter \& Gewürze!Kreuzkümmel}\\
    1 TL & Salz\index{Kräuter \& Gewürze!Salz}\\
    1{,}5 EL & Zitronensaft\index{Obst!Zitronen}
}

\preparation{%
    \step%
    Zwiebel, Knoblauch, Ingwer und Karotte klein schneiden. 
    In etwas Öl etwa 3 Minuten anbraten.
         
    \step%
    Gewürze hinzufügen und etwa 1 Minute mitrösten.\\
           
    \step%
    Linsen dazugeben und kurz (ca. 1 Minute) mit anbraten.
         
    \step%
    Passierte Tomaten, Kokosmilch und die Hälfte der Brühe hinzufügen. 
        
    \step%
    Zugedeckt etwa 20 Minuten köcheln lassen.\\
     
    \step%
    Zitronensaft und restliche Brühe dazugeben.\\
        
    \step%
    Weitere 10--15 Minuten köcheln, bis die Linsen weich, aber noch leicht bissfest sind.
         
    \step%
    Zum Schluss Margarine unterrühren.
}

\suggestion[Servieren]{%
    Am besten mit Naan-Brot oder Klebreis servieren.
}

\hint{%
    Wenn es cremiger sein soll, etwas länger köcheln lassen oder leicht pürieren.
}
    
\end{recipe}